\chapter{Clinical trials}
\index{Clinical trials}

\section*{Definition and Overview}
Clinical trials are prospective, interventional research studies involving human participants designed to evaluate the safety, efficacy, dosage, and side effects of medical, surgical, pharmacological, or behavioural interventions. They represent the gold standard of evidence-based medicine, transforming laboratory discoveries and preclinical findings into safe clinical therapies. By systematically comparing new treatments against existing standards of care or placebos, clinical trials provide the rigorous scientific data required by regulatory bodies (such as the Therapeutic Goods Administration in many countries) to approve new therapies. The primary clinical goal of any trial is to determine whether a novel intervention improves health outcomes, reduces disease burden, or enhances the quality of life of patients while maintaining an acceptable safety profile.

\section*{Epidemiology}
The volume and distribution of clinical trials globally and in many countries have expanded significantly over the past decades. Thousands of new trials are registered annually on databases such as ClinicalTrials.gov and the local New Zealand Clinical Trials Registry (ANZCTR). In many countries, clinical trials are a major sector of health and medical research, with over 1,000 new trials initiated each year, involving tens of thousands of local participants. The distribution of trials is heavily concentrated in high-burden therapeutic areas, with oncology (cancer research) representing the largest proportion, followed by cardiovascular disease, neurology, immunology, and infectious diseases. Participation rates vary by demographic factors, with adult populations being the most studied, though paediatric and geriatric trials are increasingly mandated to ensure age-specific safety and efficacy data.

\section*{Aetiology and Pathophysiology}
The scientific foundation of clinical trials relies on strict methodological principles designed to minimise bias, control confounding variables, and ensure the validity of the study's findings.
\begin{itemize}
    \item \textbf{Clinical equipoise:} The ethical requirement that there must be genuine scientific uncertainty within the medical community regarding which treatment (the new intervention or the control) is superior.
    \item \textbf{Randomisation:} The process of assigning participants to treatment groups by chance to ensure that known and unknown confounding variables are distributed equally, reducing selection bias.
    \item \textbf{Blinding:} Methodological techniques where participants (single-blind) and investigators (double-blind) are kept unaware of group assignments to prevent performance and detection bias.
    \item \textbf{Ethical oversight:} Review and approval by a Human Research Ethics Committee (HREC) to ensure that the risks to participants are minimised and reasonable in relation to anticipated benefits.
    \item \textbf{Control groups:} The use of a comparator group receiving either the current standard of care or a placebo to establish a baseline for measuring the efficacy of the investigational intervention.
\end{itemize}

\section*{Clinical Features}
Clinical trials are conducted in sequential phases, each designed to answer specific questions regarding the investigational product.
\begin{itemize}
    \item \textbf{Phase I trials:} Small-scale studies (typically 20 to 100 healthy volunteers or patients with advanced disease) designed to evaluate safety, determine safe dosage ranges, and identify pharmacological side effects.
    \item \textbf{Phase II trials:} Studies in larger patient cohorts (typically 100 to 300 individuals with the target condition) designed to assess preliminary efficacy and further evaluate safety.
    \item \textbf{Phase III trials:} Large-scale, randomised controlled studies (involving 1,000 to 3,000+ patients across multiple centres) to confirm clinical efficacy, monitor side effects, and compare the new intervention with standard treatments.
    \item \textbf{Phase IV trials:} Ongoing studies conducted after regulatory approval to monitor long-term safety, real-world effectiveness, and rare side effects in the general population.
    \item \textbf{Inclusion and exclusion criteria:} Explicit protocols defining who is eligible to participate (inclusion criteria) and who must be excluded (exclusion criteria) to ensure safety and scientific validity.
\end{itemize}

\section*{Differential Diagnosis}
Clinical trials must be distinguished from observational and non-interventional research designs, which occupy different levels on the hierarchy of scientific evidence.
\begin{itemize}
    \item \textbf{Cohort studies:} Observational studies where a group of exposed and unexposed individuals is followed over time to compare outcomes, without the researcher assigning the exposure.
    \item \textbf{Case-control studies:} Retrospective observational studies comparing individuals with a specific outcome (cases) to those without (controls) to identify risk factors.
    \item \textbf{Cross-sectional studies:} Observational studies that assess exposure and outcome simultaneously in a population at a single point in time, unable to establish temporality.
    \item \textbf{Systematic reviews and meta-analyses:} Studies that synthesise and quantitatively analyse data from multiple independent clinical trials to provide a single pooled estimate of effect.
    \item \textbf{Preclinical research:} In vitro (cell culture) and in vivo (animal model) studies conducted in laboratory settings prior to human clinical testing.
\end{itemize}

\section*{When to Seek Medical Care}
Participants enrolled in a clinical trial must be closely monitored and should contact their study coordinator or seek medical care if they experience new symptoms, worsening of their underlying condition, or suspected adverse drug reactions.

\textbf{Contact your local emergency services for:}
\begin{itemize}
    \item Severe, life-threatening symptoms such as anaphylaxis (severe allergic reaction with facial swelling and difficulty breathing).
    \item Sudden chest pain, severe shortness of breath, or cardiac arrest.
    \item Severe, uncontrollable bleeding or sudden neurological changes (e.g., loss of speech, weakness on one side).
    \item Loss of consciousness, seizures, or acute confusion.
\end{itemize}

\section*{Diagnosis and Investigations}
Before and during a clinical trial, participants undergo rigorous diagnostic testing to ensure eligibility and monitor safety parameters.
\begin{itemize}
    \item \textbf{Participant screening:} Comprehensive physical examination, medical history review, and baseline clinical tests to verify inclusion and exclusion criteria.
    \item \textbf{Laboratory safety testing:} Frequent blood chemistry, haematology (full blood count), liver and renal function tests to monitor for systemic toxicities.
    \item \textbf{Electrocardiograms and imaging:} Baseline and periodic cardiac monitoring or radiological scans to evaluate anatomical responses (e.g., tumour shrinkage in oncology trials) and detect structural toxicities.
    \item \textbf{Efficacy endpoints measurement:} Systematic tracking of primary and secondary outcomes (e.g., survival rates, blood pressure reduction) as defined by the study protocol.
    \item \textbf{Pharmacokinetic assays:} Measurements of drug absorption, distribution, metabolism, and excretion in blood or urine samples.
\end{itemize}

\section*{Management}
The conduct and management of clinical trials are strictly regulated to protect human subjects and maintain scientific integrity.
\begin{itemize}
    \item \textbf{Good Clinical Practice:} The international ethical and quality standard that governs the design, conduct, performance, monitoring, auditing, recording, analysis, and reporting of clinical trials.
    \item \textbf{Informed consent process:} A continuous process ensuring that participants are fully informed of the trial's purpose, risks, benefits, and alternatives, and that they participate voluntarily.
    \item \textbf{Data Safety Monitoring Board:} An independent committee of experts that reviews accumulating safety data to ensure participant safety and recommend trial termination if early benefit or harm is demonstrated.
    \item \textbf{Protocol adherence:} Regular audits by clinical research associates to ensure the trial is conducted in strict compliance with the approved protocol.
    \item \textbf{Pharmacovigilance:} Systematic reporting of Adverse Events (AEs) and Serious Adverse Events (SAEs) to regulatory authorities and ethics committees.
\end{itemize}

\section*{Complications}
Participating in or conducting clinical trials carries inherent medical, logistical, and ethical risks.
\begin{itemize}
    \item \textbf{Adverse drug reactions:} Unexpected or severe side effects caused by the investigational agent, ranging from mild discomfort to life-threatening events.
    \item \textbf{Protocol deviations:} Accidental or intentional departures from the trial protocol, which can compromise data quality and participant safety.
    \item \textbf{Participant attrition:} High drop-out rates due to side effects, lack of efficacy, or loss of interest, which can bias the study results (attrition bias).
    \item \textbf{Sponsor or investigator bias:} Conscious or unconscious manipulation of trial design, data analysis, or publication to favour positive outcomes.
    \item \textbf{Unintended unblinding:} Premature disclosure of treatment assignments due to distinctive side effects or analytical errors, compromising study objectivity.
\end{itemize}

\section*{Prognosis and Outcomes}
The prognosis of a clinical trial refers to its scientific success and potential to change clinical practice. The vast majority of early-phase trials do not progress to clinical approval; only about 10--15\% of drugs entering Phase I trials eventually receive regulatory clearance. For successful trials, the translation of positive findings into standard practice can take several years. The dissemination of outcomes through peer-reviewed journals and international registries is critical, regardless of whether the results are positive, negative, or neutral, to prevent publication bias and guide future medical research.

\section*{Prevention and Self-Care}
Participants in clinical trials can actively manage their safety and optimize trial validity by following key self-care and protocol guidelines.
\begin{itemize}
    \item \textbf{Adhere strictly to the study protocol:} Take study medications exactly as prescribed and attend all scheduled follow-up visits.
    \item \textbf{Maintain a detailed symptoms log:} Record any changes in physical health, side effects, or unusual symptoms in a diary to report to the study team.
    \item \textbf{Disclose concurrent treatments:} Inform the research team of any new over-the-counter medications, supplements, or prescriptions started during the trial.
    \item \textbf{Ask questions and exercise autonomy:} Seek clarification from the research staff at any point, and remember that participation is completely voluntary and consent can be withdrawn at any time.
    \item \textbf{Avoid lifestyle changes:} Maintain stable dietary and exercise habits during the trial unless instructed otherwise, to avoid introducing confounding variables.
\end{itemize}

\section*{Sources and Further Reading}
The following reputable organisations provide further information on clinical trials:
\begin{itemize}
    \item Australian Clinical Trials (Australian Government). \textit{Information on how clinical trials work and how to find trials in Australia.} https://www.clinicaltrials.gov.au/
    \item National Health and Medical Research Council (NHMRC). \textit{Guidelines on ethics and research conduct in clinical trials.} https://www.nhmrc.gov.au/
    \item ClinicalTrials.gov (US National Library of Medicine). \textit{Global registry of clinical trials and study results.} https://clinicaltrials.gov/
    \item World Health Organization (WHO). \textit{International Clinical Trials Registry Platform (ICTRP) and global trial standards.} https://www.who.int/
    \item Therapeutic Goods Administration (TGA). \textit{Information on regulation of therapeutic goods and clinical trial schemes in Australia.} https://www.tga.gov.au/
\end{itemize}
